% ═══════════════════════════════════════════════════════════════════════════
% CAPÍTULO 8 — RESULTADOS E DISCUSSÃO
% ═══════════════════════════════════════════════════════════════════════════
\chapter{Resultados e Discussão}
\label{ch:resultados}

\section{Resultados de Teste}

\subsection{Suites TDD — Evidência Completa (312/312 CTs)}

A Tabela~\ref{tab:tdd_full} apresenta o resultado consolidado das 15 suites TDD. Cada linha é verificável: o comando \texttt{python specs/test\_*.py} reproduz o resultado.

\begin{table}[H]
\centering
\small
\caption{Resultados Consolidados das Suites TDD — 312/312 PASS}
\label{tab:tdd_full}
\begin{tabular}{lcc}
\toprule
\textbf{Suite} & \textbf{SPEC} & \textbf{CTs} \\
\midrule
test\_frontmatter\_validator & 025 & 161/161 \\
test\_evolve\_pipeline & 026 & 10/10 \\
test\_evolve\_e2e & 027 & 8/8 \\
test\_noological\_scanner & 028 & 18/18 \\
test\_teleological\_scanner & 029 & 12/12 \\
test\_evolutionary\_scanner & 030 & 16/16 \\
test\_scanner\_refinement & 031 & 16/16 \\
test\_minimum\_capability\_solver & 032 & 14/14 \\
test\_capability\_composer & 033 & 13/13 \\
test\_capability\_integration & 035 & 6/6 \\
test\_metacognitive\_pipeline & 036 & 8/8 \\
test\_structural\_noise\_scanner & 037 & 8/8 \\
test\_structural\_compression\_engine & 037b & 6/6 \\
test\_n2\_n3\_upgrades & 037c & 8/8 \\
test\_behavioral\_autonomy & 038 & 8/8 \\
\midrule
\textbf{TOTAL} & & \textbf{312/312} \\
\bottomrule
\end{tabular}
\end{table}

\subsection{Evolução da Cobertura de Testes}

A progressão do número de CTs ao longo dos ciclos evolutivos demonstra uma trajetória consistente de expansão controlada. Cada novo módulo adiciona entre 6 e 22 CTs, e a exigência de zero regressão garante que o sistema nunca perde capacidades já validadas.

\subsection{Progressão dos Níveis de Consciência}

A Tabela~\ref{tab:progressao} resume a progressão dos níveis de consciência ao longo dos últimos três ciclos.

\begin{table}[H]
\centering
\small
\caption{Progressão dos Níveis de Consciência (R21-R23)}
\label{tab:progressao}
\begin{tabular}{cccl}
\toprule
\textbf{Ciclo} & \textbf{Nível} & \textbf{CTs} & \textbf{Descrição} \\
\midrule
R21 & N2.5 & 282 & N2 parcial (1/5), N3 incipiente (1/6) \\
R22 & N3.0 & 304 & N2 completo (5/5), N3 completo (4/4) \\
R23 & N3.5 & 312 & N3 completo + Behavioral Gate preventivo (8/8) \\
\bottomrule
\end{tabular}
\end{table}

\section{Estudos de Caso}

\subsection{Caso 1: Auto-Diagnóstico do Ecossistema}

\textbf{Contexto}: O pipeline de scanners foi aplicado ao próprio ecossistema OpenCode para avaliar sua capacidade de auto-diagnóstico. O corpus consistiu nos documentos AGENTS.md, SPEC\_COVERAGE.md, e SPEC-033. Os objetivos teleológicos incluíram seis metas de AGI.

\textbf{Resultados}: O scanner identificou cobertura noológica de 27\% (50-60\% real com equivalências semânticas), 4 gaps críticos e um construction cost de 6\%. Os 4 gaps foram subsequentemente implementados (SPEC-036, SPEC-037), demonstrando que o sistema não apenas diagnostica, mas também age sobre o diagnóstico.

\textbf{Interpretação}: Este caso demonstra a capacidade de auto-referência funcional do sistema — o scanner analisou o código que implementa o scanner, identificou suas próprias limitações, e as corrigiu. Este é o equivalente computacional de um paciente que se auto-diagnostica e se auto-medica.

\subsection{Caso 2: Compressão Estrutural de Texto Acadêmico}

\textbf{Contexto}: O Structural Compression Engine (SCE) foi aplicado a um corpus de 300 palavras contendo estruturas argumentativas, exemplos e ruído.

\textbf{Resultados}: CR de 2.2x, CPS de 100\%, FLI de 0\%, DG de 2.2.

\textbf{Interpretação}: O SCE reduziu o volume de texto pela metade sem perda de estrutura cognitiva. A preservação de 100\% (CPS=1.0) indica que o modelo reduzido capturou todas as funções essenciais do original.

\subsection{Caso 3: Behavioral Gate em Produção}

\textbf{Contexto}: O TrustEngine foi testado em cenário simulado de produção com 3 tipos de ações.

\textbf{Resultados}: Ação confiável classificada como ``safe'' (trust 1.00), ação arriscada como ``risky'' (trust 0.32), e ação nova operando em shadow mode (trust $\leq$ 0.50).

\textbf{Interpretação}: O BehavioralGate demonstrou capacidade de distinguir entre ações seguras e arriscadas baseado em histórico de outcomes, e de limitar a confiança em ações novas (shadow mode). Esta é a base para prevenção de falhas em sistemas autônomos.

\section{Comparação com o Estado da Arte}

A Tabela~\ref{tab:comparacao_final} compara o OpenCode com sistemas relacionados.

\begin{table}[H]
\centering
\small
\caption{Comparação com Sistemas Relacionados}
\label{tab:comparacao_final}
\begin{tabular}{lccccc}
\toprule
\textbf{Característica} & \textbf{OpenCode} & \textbf{NEXO} & \textbf{Ruflo} & \textbf{AutoGPT} & \textbf{LangChain} \\
\midrule
Auto-observação (N2) & Sim & Parcial & Não & Não & Não \\
Detecção anomalias (N3) & Sim & Parcial & Não & Não & Não \\
Correção automática (N3) & Sim & Não & Não & Não & Não \\
Inferência causal & Sim & Não & Não & Não & Não \\
Gate preventivo (N3.5) & Sim & Sim & Não & Não & Não \\
Governança Ostrom & Sim & Não & Não & Não & Não \\
Compressão estrutural & Sim & Não & Não & Não & Não \\
TDD (312 CTs) & Sim & Não & Não & Parcial & Não \\
Ciclos evolutivos & 23 & — & — & — & — \\
\bottomrule
\end{tabular}
\end{table}

\section{Limitações (Transparência Total)}

Oito limitações são reconhecidas:

\begin{enumerate}
    \item \textbf{Decomposição limitada}: Apenas 10 das 92 categorias do NoologicalScanner têm templates de decomposição. As demais caem em frontier (cost=1.0).
    \item \textbf{Biblioteca estática}: Os 85 insumos cognitivos são estáticos; não há mecanismo de auto-expansão da biblioteca.
    \item \textbf{Auto-representação simbólica}: O SelfModel N0-N3 opera sobre métricas declarativas, sem componente experiencial ou neural.
    \item \textbf{Aprendizado heurístico}: Thresholds são ajustados por heurísticas (falsos positivos), não por gradient descent ou otimização formal.
    \item \textbf{Semântica ausente}: O sistema opera sobre palavras (keyword matching), não sobre significados (embeddings, grafos semânticos).
    \item \textbf{Intencionalidade externa}: Goals são fornecidos como strings pelo operador; o sistema não gera objetivos próprios.
    \item \textbf{SPS limitado}: O Structural Preservation Score fica abaixo de 0.90 em textos muito curtos ($<$ 10 sentenças).
    \item \textbf{Trust não-transferível}: O Trust Scoring não propaga confiança entre ações similares (ex.: se confio em ``scan\_raciocinio'', não confio automaticamente em ``scan\_metodos'').
\end{enumerate}

\section{Trabalhos Futuros}

\begin{enumerate}
    \item \textbf{SPEC-034}: Motor de decomposição avançado com analogia polimática e LLM para categorias sem template.
    \item \textbf{SPEC-039}: Aprendizado contínuo com gradient-based threshold optimization.
    \item \textbf{SPEC-040}: Compreensão semântica via embeddings e grafos de conhecimento dinâmicos.
    \item \textbf{SPEC-041}: Intencionalidade emergente via homeostase computacional.
    \item \textbf{SPEC-042}: Cross-action trust propagation para transferência de confiança entre ações similares.
\end{enumerate}
